%%=============================================================================
%% Methodologie
%%=============================================================================

\chapter{\IfLanguageName{dutch}{Methodologie}{Methodology}}%
\label{ch:methodologie}

%% TODO: In dit hoofstuk geef je een korte toelichting over hoe je te werk bent
%% gegaan. Verdeel je onderzoek in grote fasen, en licht in elke fase toe wat
%% de doelstelling was, welke deliverables daar uit gekomen zijn, en welke
%% onderzoeksmethoden je daarbij toegepast hebt. Verantwoord waarom je
%% op deze manier te werk gegaan bent.
%% 
%% Voorbeelden van zulke fasen zijn: literatuurstudie, opstellen van een
%% requirements-analyse, opstellen long-list (bij vergelijkende studie),
%% selectie van geschikte tools (bij vergelijkende studie, "short-list"),
%% opzetten testopstelling/PoC, uitvoeren testen en verzamelen
%% van resultaten, analyse van resultaten, ...
%%
%% !!!!! LET OP !!!!!
%%
%% Het is uitdrukkelijk NIET de bedoeling dat je het grootste deel van de corpus
%% van je bachelorproef in dit hoofstuk verwerkt! Dit hoofdstuk is eerder een
%% kort overzicht van je plan van aanpak.
%%
%% Maak voor elke fase (behalve het literatuuronderzoek) een NIEUW HOOFDSTUK aan
%% en geef het een gepaste titel.

Zoals eerder aangehaald werd, richt dit onderzoek zich op het vinden van een oplossing die commerciële dark patterns van in Vlaanderen vaak gebruikte streamingdiensten automatisch opspoort en aan de gebruikers signaleert. Om de onderzoeksvraag te beantwoorden, werd een weloverwogen plan gevolgd dat bestaat uit \textbf{vijf verschillende fasen}. In eerste instantie werd de huidige stand van zaken onderzocht aan de hand van een literatuurstudie en werd een longlist opgesteld van alle bestaande dark pattern detectietools. Nadien werden uit het literatuuronderzoek de requirements waaraan de tool moet voldoen opgesteld. Aan de hand van deze vereisten werd de longlist uit fase 1 vernauwd tot een shortlist die de meest geschikte tools bevat voor deze casus. Hieruit werd de meest beloftevolle bestaande open-source tool gekozen voor deze casus om verder mee aan de slag te gaan tijdens de uitwerking van een \acrlong{poc} om zo de tool te optimaliseren voor dark pattern detectie op streamingdiensten. In de laatste fase werd een conclusie geformuleerd. 

\vspace{1em}

\section{Stand van zaken \& Longlist}
\label{sec:stand-van-zaken-en-longlist}

\vspace{1em}

In de eerste fase van het onderzoek werd het onderzoeksdomein uitvoerig verkend om op die manier een goed theoretisch kader op te bouwen. Via literatuuronderzoek werd onderzocht wat dark patterns zijn, welke er bestaan, welke negatieve gevolgen deze kunnen hebben en hoe bewust mensen zijn van de patronen en hun gevolgen. Verder werd nagegaan waar deze patronen het vaakst voorkomen (websites en apps) en of er verschillen zijn per platform. Daarnaast werd ook de huidige wet- en regelgeving en de link tussen businessdoelstellingen van streamingdiensten en het gebruik van dark patterns onder de loep genomen. Aanvullend aan dit literatuuronderzoek, werd in kaart gebracht welke streamingdiensten het populairst zijn in Vlaanderen en welke dark patterns zij toepassen. De literaire studie werd gedaan gebruikmakend van hoofdzakelijk \emph{Google Scholar} en wetenschappelijke databanken zoals \emph{ACM digital library} en \emph{Springer Online Journals}.

\vspace{1em}

Op basis van literatuuronderzoek werd er een \textbf{longlist} opgesteld van digitale oplossingen die gebruikers kunnen helpen bij het herkennen van dark patterns op streamingdiensten. Ieder alternatief werd gedocumenteerd met voor- en nadelen op basis van voorafbepaalde aspecten. Het doel van de longlist was om een zo compleet mogelijk overzicht te creëren van bestaande tools, dat in een verdere fase gebruikt werd voor het opstellen van de shortlist. 

\vspace{1em}

De inzichten uit deze fase vormden een belangrijke basis voor de verdere fasen van de bachelorproef, waaronder de uitvoering van de requirementsanalyse.
Deze eerste fase resulteerde in een opsomming van dark patterns die voorkomen bij populaire streamingdiensten, een overzicht van bestaande Europese-wetgeving rond dark patterns, een oplijsting van de belangrijkste bevindingen uit de literatuurstudie en een longlist van reeds ontwikkelde detectietools, inclusief hun voor- en nadelen.

\vspace{1em}

Door deze eerste fase af te ronden werd voldoende informatie verzameld om de volgende deelvragen te beantwoorden:

\begin{itemize}
    \item Welke dark patterns komen het meest voor op websites en apps van populaire streamingdiensten die door Vlamingen tussen 18 en 34 jaar gebruikt worden?
    \item In welke mate zijn gebruikers zich bewust van de aanwezigheid van dark patterns bij het gebruik van streamingdiensten?
    \item Welke verbanden bestaan er tussen de aanwezigheid van specifieke dark patterns en de businessdoelstellingen van streamingdiensten?
    \item Op welk platform (website of mobiele app) komen dark patterns het meest frequent voor?
    \item Kunnen technologische tools helpen bij het verhogen van bewustzijn rond bepaalde onderwerpen?
    \item Welke wet- en regelgeving bestaat er rond het gebruik van dark patterns in België en de Europese Unie?
    \item Welke bestaande tools die gebruikers waarschuwen voor dark patterns bestaan er al, en wat zijn hun sterke en zwakke punten?
    \item Op welke manier kunnen dark patterns aan gebruikers gesignaleerd worden zonder dat hun gebruikerservaring negatief beïnvloed wordt?
\end{itemize}

\vspace{1em}

\section{Requirements verzamelen}
\label{sec:requirements-verzamelen}

\vspace{1em}

In fase~2 ligt de focus op het verzamelen van alle vereisten waaraan de nieuwe tool moet voldoen. Op basis van de resultaten uit het literatuuronderzoek werden de requirements opgesteld. Er werd onder andere rekening gehouden met welke dark patterns het belangrijkste zijn, waar de tool moet ingezet worden, welke aspecten overgenomen kunnen worden van andere tools en welke toevoegingen nodig zijn voor een betere digitale oplossing. Aanvullend werden de voorlopige requirements afgetoetst en aangevuld via interviews met een UI/UX-designer, front-end developer en twee personen uit de doelgroep. Het einde van deze fase werd gekenmerkt door het opstellen van een volledige en geprioriteerde lijst van criteria die het \acrlong{mvp} (\acrshort{mvp}) definiëren.

\vspace{1em}

Het concrete resultaat uit deze fase bestaat uit een lijst van functionele en niet-functionele requirements, geordend volgens hun belang. Voor prioritering werd de \emph{MoSCoW}-methode gebruikt.

\vspace{1em}

Door het afronden van de tweede fase werd voldoende informatie verzameld om de volgende deelvraag deels te beantwoorden:

\begin{itemize}
    \item Welke digitale of interactieve oplossing is geschikt voor het automatisch detecteren en signaleren van dark patterns op streamingdiensten, en welke (niet-)functionele vereisten zijn hiervoor essentieel?
\end{itemize}

\vspace{1em}

\section{Shortlist}
\label{sec:shortlist}

\vspace{1em}

Vooraleer een \acrlong{poc} uitgewerkt kon worden, werd er gekeken welke oplossing uit de longlist het meest beloftevol en waardevol is voor het onderzoek. Daarvoor werden in een derde fase selectiecriteria opgesteld op basis van de inzichten uit de literatuurstudie en de requirementsanalyse uit fase~2. Nadien werden deze criteria toegepast in twee opeenvolgende selectierondes door middel van overzichtelijke vergelijkingstabellen.

\vspace{1em}

Op basis van deze selectierondes werd een shortlist opgesteld, waaruit de meest beloftevolle oplossing geselecteerd werd door middel van een praktische evaluatie. Tot slot werd de keuze voor deze tool grondig onderbouwd door middel van een duidelijke motivatie. Met de gekozen tool werd in de voorlaatste fase verder aan de slag gegaan.

\vspace{1em}

Door het afronden van de derde fase werd voldoende informatie verzameld om de volgende deelvraag volledig te beantwoorden:

\begin{itemize}
    \item Welke digitale of interactieve oplossing is geschikt voor het automatisch detecteren en signaleren van dark patterns op streamingdiensten, en welke (niet-)functionele vereisten zijn hiervoor essentieel?
\end{itemize}

\section{Proof of Concept}
\label{sec:poc}

\vspace{1em}

Fase~4 omvat de uitwerking van een \acrlong{poc}. De geselecteerde oplossing uit fase~3 werd in deze fase verder uitgewerkt tot een eerste \acrlong{mvp} door de tool aan te passen en uit te breiden voor de specifieke use case van dit onderzoek.

\vspace{1em}

In eerste instantie werden er mockups uitgewerkt in \emph{Figma} om de algemene flow van de tool te bepalen. Nadien werden in overleg met de copromotoren en aan de hand van een korte aanvullende literatuurstudie de meest geschikte technologieën (zoals \emph{Kotlin}, \emph{Python}, \ldots) gekozen. Daarna startte de implementatie van de tool.

\vspace{1em}

Deze fase resulteerde in een werkend prototype die de belangrijkste functionaliteiten van de tool bevat, aangevuld met een korte documentatie over de werking. Eveneens werd gevalideerd of de uitgewerkte tool al dan niet een oplossing biedt voor het gestelde probleem. Dit werd gedaan aan de hand van een kleine test bij enkele testgebruikers die tot de doelgroep behoren. Zij voerden een aantal taken uit op een streamingplatform dat gebruikmaakt van dark patterns, zowel zonder als met de tool. Uit deze test werd een eerste indicatie afgeleid over hoe gebruiksvriendelijk de tool is door te polsen naar hun gebruikerservaring met de tool, in welke mate ze de patronen herkennen en in hoe snel zij deze waarnemen. De resultaten hiervan gaven een eerste indicatie in de mogelijke meerwaarde van de \acrshort{poc}, zonder ervan uit te gaan dat dit geldt voor de hele doelgroep.

\vspace{1em}

Tot slot werd nagegaan in welke mate de gerealiseerde \acrshort{poc} voldoet aan de requirements die in fase~2 werden opgesteld. Hierbij werd per requirement nagegaan of deze volledig, gedeeltelijk of niet werd gerealiseerd. Naast de gerealiseerde functionaliteiten werden ook beperkingen en openstaande requirements geïdentificeerd. Deze evaluatie vormde samen met de gebruikerstest de basis voor het beoordelen van de werking van de uitgewerkte oplossing en het formuleren van mogelijke verbeterpunten.

\vspace{1em}

Deze evaluatie maakte het mogelijk om de laatste deelvraag van dit onderzoek te beantwoorden:

\begin{itemize}
    \item In welke mate voldoet de uitgewerkte digitale oplossing aan de opgestelde requirements en welke verbeterpunten kunnen worden geïdentificeerd?
\end{itemize}

\vspace{1em}

\section{Conclusie}
\label{sec:conclusie}

\vspace{1em}

In de laatste fase van het onderzoek werden de resultaten van het onderzoek samengebracht en geïnterpreteerd. Op basis van de uitgewerkte \acrshort{poc}, de evaluatie van de requirements en de resultaten van de gebruikerstesten werd een antwoord geformuleerd op de verschillende deelvragen en de centrale onderzoeksvraag.

\vspace{1em}

Daarnaast werd gereflecteerd over de gerealiseerde oplossing. Hierbij werd besproken in welke mate de \acrshort{poc} een meerwaarde biedt voor het herkennen van dark patterns op streamingdiensten en welke beperkingen tijdens het onderzoek naar voren zijn gekomen. Ook werd nagegaan welke aspecten van de oplossing nog verder ontwikkeld of onderzocht moeten worden.

\vspace{1em}

Tot slot werden mogelijke richtingen voor verder onderzoek besproken. Hierbij wordt onder andere aandacht besteed aan aspecten die buiten de scope van dit onderzoek vielen, maar tijdens de ontwikkeling of evaluatie relevant bleken.
